\documentclass[11pt]{article}

\usepackage[margin=1in]{geometry}
\usepackage{amsmath}
\usepackage{amssymb}
\usepackage{siunitx}
\usepackage{booktabs}
\usepackage{enumitem}
\usepackage[hidelinks]{hyperref}

\title{Low Cycle Fatigue: Physics and Equations for Review}
\author{lcf-strain-life}
\date{\today}

\begin{document}
\maketitle

\begin{abstract}
This is a self-contained statement of the fatigue physics implemented in the
toolkit, written for review by a materials or fatigue specialist. It contains the
equations, the symbols and units, the conventions, and the source for each
relation. It contains no software detail. A short sign-off table is at the end.
\end{abstract}

\section{Conventions}

All analysis uses true stress and true strain. Stress and modulus are in
\si{\mega\pascal}. Strain is a dimensionless fraction. Life is expressed in
reversals \(2N_f\), where one cycle is two reversals. The fatigue strength
exponent \(b\) and the fatigue ductility exponent \(c\) are negative.

\begin{table}[h]
\centering
\small
\begin{tabular}{lll}
\toprule
Symbol & Meaning & Unit \\
\midrule
\(\sigma,\ \varepsilon\) & true stress, true strain & \si{\mega\pascal}, - \\
\(\Delta\sigma/2,\ \Delta\varepsilon/2\) & stress, total strain amplitude & \si{\mega\pascal}, - \\
\(\Delta\varepsilon_e/2,\ \Delta\varepsilon_p/2\) & elastic, plastic strain amplitude & - \\
\(\sigma_m,\ \sigma_{max}\) & mean stress, maximum stress & \si{\mega\pascal} \\
\(E\) & Young's modulus & \si{\mega\pascal} \\
\(\sigma'_f,\ b\) & fatigue strength coefficient, exponent & \si{\mega\pascal}, - \\
\(\varepsilon'_f,\ c\) & fatigue ductility coefficient, exponent & - \\
\(K',\ n'\) & cyclic strength coefficient, strain-hardening exponent & \si{\mega\pascal}, - \\
\(2N_f\) & reversals to failure & - \\
\(K_t,\ K_f,\ q\) & stress concentration, fatigue notch factor, notch sensitivity & - \\
\bottomrule
\end{tabular}
\end{table}

\section{True stress and strain}
Engineering values \(e\) and \(\sigma_\text{eng}=F/A_0\) convert to true values by
\begin{equation}
\varepsilon = \ln(1+e), \qquad \sigma = \sigma_\text{eng}(1+e).
\end{equation}
Valid up to necking. Reference: Dowling, Mechanical Behavior of Materials, 4th ed.

\section{Cyclic stress-strain, Ramberg-Osgood}
\begin{equation}
\varepsilon = \frac{\sigma}{E} + \left(\frac{\sigma}{K'}\right)^{1/n'},
\qquad
\Delta\varepsilon = \frac{\Delta\sigma}{E} + 2\left(\frac{\Delta\sigma}{2K'}\right)^{1/n'}.
\end{equation}
The second form is the doubled hysteresis branch. Reference: Ramberg and Osgood
1943, NACA TN 902. Dowling 4th ed., Eq. 14.12.

\section{Strain-life}
Elastic, Basquin:
\begin{equation}
\frac{\Delta\sigma}{2} = \sigma'_f (2N_f)^{b},
\qquad
\frac{\Delta\varepsilon_e}{2} = \frac{\sigma'_f}{E}(2N_f)^{b}.
\end{equation}
Plastic, Coffin-Manson:
\begin{equation}
\frac{\Delta\varepsilon_p}{2} = \varepsilon'_f (2N_f)^{c}.
\end{equation}
Total strain-life and the elastic-plastic transition life:
\begin{equation}
\frac{\Delta\varepsilon}{2} = \frac{\sigma'_f}{E}(2N_f)^{b} + \varepsilon'_f (2N_f)^{c},
\qquad
2N_t = \left(\frac{\varepsilon'_f E}{\sigma'_f}\right)^{1/(b-c)}.
\end{equation}
Compatibility, checked and reported but not forced:
\begin{equation}
n' = \frac{b}{c}, \qquad K' = \frac{\sigma'_f}{(\varepsilon'_f)^{b/c}}.
\end{equation}
The plastic line is fit over the low cycle regime, excluding near-runout points
whose plastic strain is at measurement noise level. References: Basquin 1910,
Coffin 1954, Manson 1953, Dowling 4th ed. Eq. 14.3 to 14.6.

\section{Mean-stress corrections}
Morrow, elastic term shifted by the mean stress:
\begin{equation}
\frac{\Delta\varepsilon}{2} = \frac{\sigma'_f-\sigma_m}{E}(2N_f)^{b} + \varepsilon'_f (2N_f)^{c}.
\end{equation}
Modified Morrow, both terms shifted:
\begin{equation}
\frac{\Delta\varepsilon}{2} = \frac{\sigma'_f-\sigma_m}{E}(2N_f)^{b}
+ \varepsilon'_f\left(\frac{\sigma'_f-\sigma_m}{\sigma'_f}\right)^{c/b}(2N_f)^{c}.
\end{equation}
Smith-Watson-Topper:
\begin{equation}
\sigma_{max}\,\frac{\Delta\varepsilon}{2} = \frac{(\sigma'_f)^2}{E}(2N_f)^{2b}
+ \sigma'_f \varepsilon'_f (2N_f)^{b+c},
\qquad
\sigma_{ar} = \sqrt{\sigma_{max}\,\sigma_a}.
\end{equation}
Walker, with the equivalent fully-reversed amplitude and the steel estimate of
the exponent:
\begin{equation}
\sigma_{ar} = \sigma_{max}^{\,1-\gamma}\,\sigma_a^{\,\gamma},
\qquad
\gamma_{steel} = 0.8818 - 2.00\times10^{-4}\,\sigma_u.
\end{equation}
Morrow equivalent fully-reversed amplitude:
\begin{equation}
\sigma_{ar} = \frac{\sigma_a}{1 - \sigma_m/\sigma'_f}.
\end{equation}
With \(\gamma=0.5\) Walker reduces to SWT. References: Morrow 1968. Smith, Watson,
Topper 1970, J. Materials 5(4):767. Walker 1970. Dowling, Calhoun, Arcari 2009,
Fatigue Fract. Eng. Mater. Struct. (steel \(\gamma\)). Dowling 4th ed. Eq. 9.18 to 9.21.

\section{Hysteresis energy and cyclic response}
The plastic strain energy density per cycle is the closed loop area
\begin{equation}
W = \oint \sigma\,d\varepsilon.
\end{equation}
The tension-compression asymmetry of a cycle is
\(R_{TC} = \lvert\sigma_{max}\rvert / \lvert\sigma_{min}\rvert\). Peak and valley
stress versus cycle give the cyclic hardening and softening response.

\section{Variable amplitude, rainflow counting}
Irregular histories are reduced to closed cycles by the rainflow method, which
pairs reversals into hysteresis loops while preserving their order in time.
Reference: ASTM E1049-85(2017), Downing and Socie 1982, Matsuishi and Endo 1968.

\section{Cumulative damage}
Palmgren-Miner linear damage, failure at a critical sum, the default being one:
\begin{equation}
D = \sum_i \frac{n_i}{N_{f,i}}.
\end{equation}
Double Linear Damage Rule, Manson-Halford, with the Phase I life fraction
referenced to the longest life in the spectrum:
\begin{equation}
f_I = 0.35\left(\frac{N_f}{N_{long}}\right)^{0.25},
\qquad N_I = N_f f_I, \quad N_{II} = N_f(1-f_I).
\end{equation}
Phase I accumulates to one, then Phase II accumulates to one. Corten-Dolan, where
\(\sigma_1\) is the maximum stress in the block and \(\alpha_i\) the cycle fractions:
\begin{equation}
N = \frac{N_{f,1}}{\sum_i \alpha_i (\sigma_i/\sigma_1)^{d}}.
\end{equation}
References: Palmgren 1924, Miner 1945, Manson and Halford 1981 (Int. J. Fracture
17:169), Corten and Dolan 1956.

\section{Notch local-strain approach}
Neuber, combined with the cyclic curve, and its range form:
\begin{equation}
\frac{(K_t S)^2}{E} = \sigma\,\varepsilon,
\qquad
\frac{(K_t \Delta S)^2}{E} = \Delta\sigma\,\Delta\varepsilon.
\end{equation}
Glinka equivalent strain energy density:
\begin{equation}
\frac{(K_t S)^2}{2E} = \frac{\sigma^2}{2E} + \frac{\sigma}{n'+1}\left(\frac{\sigma}{K'}\right)^{1/n'}.
\end{equation}
Fatigue notch factor and notch sensitivity:
\begin{equation}
K_f^{\text{Peterson}} = 1 + \frac{K_t-1}{1+a/r},
\quad
K_f^{\text{Neuber}} = 1 + \frac{K_t-1}{1+\sqrt{\beta/r}},
\quad
q = \frac{K_f-1}{K_t-1}.
\end{equation}
Neuber tends to overestimate and Glinka to underestimate the local strain. The
measured value usually lies between them. References: Neuber 1961, Molski and
Glinka 1981, Peterson 1974.

\section{Statistics and design curves}
Life is the dependent variable in the linearized regression,
\(\log_{10} N = A + B \log_{10}(\Delta\varepsilon/2)\). Confidence and prediction
intervals use the residual standard error and the Student \(t\) quantile. A
reliability and confidence design curve is the mean reduced by \(k\,s\), where
\(k\) is the one-sided Owen tolerance factor. Right-censored runouts are handled
by maximum likelihood rather than deletion. References: ASTM E739, withdrawn 2024
and used as the de facto reference, Owen 1963, Williams, Lee and Rilly 2003 (Int.
J. Fatigue 25:427).

\section{Elevated temperature}
Frequency-modified Coffin-Manson, coefficient form:
\begin{equation}
C_f = C_o\left(\frac{f}{f_{ref}}\right)^{k-1},
\qquad
\frac{\Delta\varepsilon_p}{2} = C_f (2N_f)^{c}.
\end{equation}
Linear time-fraction creep-fatigue damage, fatigue plus creep:
\begin{equation}
D = \sum_i \frac{n_i}{N_{f,i}} + \sum_j \frac{t_j}{t_{r,j}},
\end{equation}
checked against a bilinear creep-fatigue interaction envelope. References: Coffin
1971, Robinson 1952, ASTM E2714-13(2020).

\section{Multiaxial, survey only}
Critical-plane parameters, provided for evaluation. The full plane search is not
yet implemented.
\begin{equation}
P_{FS} = \frac{\Delta\gamma_{max}}{2}\left(1 + k\frac{\sigma_{n,max}}{\sigma_y}\right),
\quad
P_{BM} = \frac{\Delta\gamma_{max}}{2} + S\,\Delta\varepsilon_n,
\quad
P_{SWT} = \sigma_{n,max}\frac{\Delta\varepsilon_1}{2}.
\end{equation}
References: Fatemi and Socie 1988 (Fatigue Fract. Eng. Mater. Struct. 11(3):149),
Brown and Miller 1973, Smith, Watson, Topper 1970.

\section{Conventions and assumptions to confirm}
\begin{itemize}[nosep]
\item True stress-strain conversion assumed valid up to necking.
\item Plastic strain amplitude taken as \(\Delta\varepsilon_t/2 - \Delta\sigma/(2E)\).
\item Failure criterion is a percent load drop from the stabilized half-life peak load, default 30 percent.
\item Default mean-stress model for variable amplitude is SWT, notch default is Neuber, damage default is Miner.
\item The plastic strain-life line is fit over the low cycle regime only.
\end{itemize}

\section{Reviewer sign-off}
Reviewer name and date, then a verdict per section. Suggested verdicts: OK, OK
with note, or needs change.

\begin{table}[h]
\centering
\small
\begin{tabular}{p{6.5cm} p{2.5cm} p{5cm}}
\toprule
Section & Verdict & Note \\
\midrule
True stress and strain & & \\
Ramberg-Osgood & & \\
Strain-life & & \\
Mean-stress corrections & & \\
Energy and cyclic response & & \\
Rainflow counting & & \\
Cumulative damage & & \\
Notch local-strain & & \\
Statistics and design curves & & \\
Elevated temperature & & \\
Multiaxial survey & & \\
Conventions and assumptions & & \\
\bottomrule
\end{tabular}
\end{table}

\end{document}
